\chapter{Glossary}\label{app:glossary}\label{ch:glossary}

\begin{goals}
\textbf{In this appendix}: the words micro uses in a particular sense, each with a short definition
and the chapter that explains it. Geological terms are listed only where micro means something
specific by them.
\end{goals}

\begin{description}

\item[accept] To put your name to a finished run, making it the baseline for later runs of the same
  routine. Recorded in \file{accepted.json} beside the run record, which is never edited.
  \cref{ch:checks}.

\item[accumulation] Grade × length: what a grade–tonnage curve on a drillhole layer (and the
  \code{xtab} tool run on one) gives in place of metal, in units such as \%·m or g/t·m (times density when one is given).
  \cref{sec:gt-holes,ch:resources}.

\item[active layer] The layer micro's own panels and single keys act on: the one the filter box
  shows, \key{t} opens the table of, and \key{h} hides. Step it with \key{[} and \key{]}.
  \cref{ch:start}.

\item[atlas] A set of pages made from one composer figure, one page per bench or per section,
  written by a routine. \cref{ch:atlas}.

\item[attitude] The orientation of a plane fitted to a painted patch of an outcrop mesh, given as
  dip and dip direction, with a measure of how well the patch fits a plane. \cref{ch:structure}.

\item[attribute table] A layer's records as rows in a floating window, read in windows so a model
  of millions of rows scrolls at once. Linked to the view both ways: click a row to pick the
  record. \cref{ch:columns}.

\item[band] One slice of a swath: the records whose position along the swath direction falls in one
  interval. Each band gives one point of the swath plot. \cref{ch:swath}.

\item[band ghost] The slab of a swath band drawn in the 3D view while you hover its point on the
  plot. \cref{ch:swath}.

\item[baseline] The earlier run of a routine that a new run is judged against: the most recently
  accepted run or, if none has been accepted, the latest that finished ok. Checks read it as
  \code{prev.}. \cref{ch:checks}.

\item[blank] A missing value. Never a zero: a blank grade is unestimated or unassayed, and it is
  left out of means, curves and color scales. \cref{ch:columns,app:expr}.

\item[block edges] The outline drawn around every block, and around every cut block on a section.
  \key{⇧B} toggles it. \cref{ch:color}.

\item[bookmark] A saved viewpoint: the camera and the section. A numbered bookmark is a view
  slot. \cref{ch:sections}.

\item[broadcast] Landing a collar column on the intervals of the same holes, matched by hole id, so
  that interval-level work can use it. \cref{ch:drillholes}.

\item[calculated column (ƒ)] A column defined by an expression and computed when needed; what is
  stored is the definition. Shown with an ƒ. Materializing it makes a stored column.
  \cref{ch:filter}.

\item[captions: in the section] The default setting of a composer view for the statistics its
  captions quote (\code{\{count\}}, \code{\{mean:…\}}): only the records inside the view's section
  are counted, so each bench of an atlas reports its own numbers and the benches add up to the
  model. The other setting, \emph{whole model}, counts every record the filters keep.
  \cref{sec:compo-statsscope,ch:atlas}.

\item[category] The text column a layer is classed by (a lithology, a domain), with one legend
  entry, color and eye per value. \cref{ch:color}.

\item[check] A step of a routine that judges a result: one condition over the values of earlier
  steps, or over every row of a step's table. It passes, fails, or is not evaluated.
  \cref{ch:checks}.

\item[class] One interval of a classed numeric color scale (the values between two breaks), or one
  entry of a category legend. Each has its own color and eye. \cref{ch:color}.

\item[collar] The record of a hole's start: its id, its position and usually its length. In micro,
  the collars are one location of a drillhole set, with a table of their own. \cref{ch:drillholes}.

\item[command palette] The searchable list of every action, filtered to what applies right now.
  \key{⌘K}. \cref{ch:start}.

\item[compare] A routine step that sets a table of this run beside the same table of the baseline,
  row by row on the key, with the differences. Given a tolerance, it is also a check. \cref{ch:checks}.

\item[compatible grids] Two block models or grids whose lattices can be brought onto a common one:
  the same orientation, with pitches and origins that line up. A spatial join or a reconcile needs
  them. \cref{ch:joins}.

\item[composer] The page layout window: frames holding 3D views, charts, tables, text and legends,
  exported as a PDF that carries its own figure. \cref{ch:composer}.

\item[composite] One interval table of a drillhole set re-cut into equal lengths, numbers weighted
  by length. The result is a new points layer, each composite at its centroid.
  \cref{ch:drillholes}.

\item[contact plot] Mean grade in bins of distance from a domain boundary, on both sides of it:
  down the holes from every A-to-B transition (\ui{Contact plot…}), or by signed distance to a
  wireframe (\ui{Distance to a surface…}). The routine tool is \code{contact}. \cref{ch:contact}.

\item[cut mode] How one layer responds to the section: follow it, follow with its own slab, keep
  only the front or the back, or ignore it. \cref{ch:sections}.

\item[declared lattice] A lattice that a file or a host states instead of leaving micro to find
  it. Applied without asking, and checked against every record. \cref{ch:blocks,app:host}.

\item[declustering] Weighting samples so that densely drilled areas do not dominate a mean. micro
  uses it when it compares a model with its samples. \cref{ch:validation}.

\item[decorations] The scale bar, north arrow, color legend and title drawn on the view, live and
  in an exported figure. \cref{ch:figures}.

\item[derived column] Any column that did not come from the source file: calculated, stored,
  painted, joined or broadcast. All behave like source columns. \cref{ch:filter}.

\item[desurvey] Computing where each interval lies in space from the collar position and the survey
  stations. The method and dip convention are settings of the drillhole set. \cref{ch:drillholes}.

\item[detection limit] A lab's lower or upper reporting bound, written in assay files as
  \code{<0.01} or \code{>10}. micro reads such a column as numeric and substitutes by a rule you
  choose. \cref{ch:columns}.

\item[discovery] The pass micro makes over a new block-model file to find its lattice, extents and
  categories. Its result is cached, and a file that declares its lattice skips it.
  \cref{ch:blocks}.

\item[dogleg] How sharply a hole bends between two survey stations, in degrees per 30 m. Reported
  by the drillhole QA. \cref{ch:qa}.

\item[domaining] Assigning records to domains: by painting, by flagging against solids, or by
  \ui{assign domains}, which gives each block to the solid holding most of it.
  \cref{ch:picking,ch:geometry}.

\item[drape] Coloring a relief surface by the values of another grid, sampled at each position.
  \cref{ch:grids}.

\item[drawer] The panel under the filter box that shows the current expression as sliders, chips
  and groups. Every control rewrites its own part of the text. \cref{ch:filter}.

\item[drillhole set] A collar table, a survey table and one or more interval tables that belong
  together. In the layers panel it is one group. \cref{ch:drillholes}.

\item[element] One dataset as micro models it: a geometry, one or more locations of records, their
  columns, and the recorded operations that derived anything. \cref{ch:micro}.

\item[element manifest] The readable \file{.element.json} file beside a data file in a project,
  describing the element and listing its operations. A drillhole set has one for the whole set,
  the set manifest. \cref{ch:store}.

\item[exact / streamed] The two ways a large mesh is held: exact keeps every triangle (needed for
  solid queries and mesh QA); streamed is for display only and light on memory. \cref{ch:meshes}.

\item[exclusive group] A group in which at most one child shows at a time. \key{\{} and \key{\}}
  step through it. \cref{ch:start}.

\item[exploded view] A group's layers pulled apart along an axis so that each can be seen whole.
  Display only: picks, filters, analyses and exports read the true coordinates.
  \cref{ch:sections}.

\item[expression] A condition or a formula over the columns of one record, in micro's one
  expression language. \cref{app:expr}.

\item[F1, F2, F3] The reconciliation factors: F1 = grade control ÷ resource model, F2 = mill ÷
  grade control, F3 = mill ÷ resource model = F1 × F2, each for tonnes, grade and metal, inside the
  volume mined in a period. \cref{ch:factors}.

\item[filter] An expression that decides which records of a layer take part: the rest leave the
  view, the table and every analysis. \cref{ch:filter}.

\item[fine lattice] The lattice of a sub-blocked model at its smallest block size, on which every
  sub-block's position is expressed. \cref{ch:blocks}.

\item[fingerprint] A hash that identifies what a result was made from. A source file's is taken
  from samples across the file; a routine step's covers its inputs and settings, and decides
  whether an earlier result can be reused. \cref{ch:records,ch:store}.

\item[fraction inside] A number column holding, for each block, the share of its volume inside a
  solid, or below or above a surface, from 0 to 1. The record panel and the export dialog mark it
  \%, the one stored column not marked ✎. \cref{sec:geometry-do}.

\item[frame] A rectangle on a composer page holding one thing: a 3D view, a chart, a table, text, a
  legend, an image. \cref{ch:composer}. (In the host API, also the coordinate frame: the CRS a
  layer's coordinates are in.)

\item[global bias] The difference between a model's mean grade and the declustered mean of its
  samples, best measured inside a domain the holes sampled. \cref{ch:validation,sec:validation-domain}.

\item[grade shell] The surface around every block at or above a cutoff, made as a watertight mesh
  layer. \cref{ch:blocks}.

\item[grade–tonnage curve] Tonnage and mean grade above each of a list of cutoffs, weighted by
  block volume and density (volume alone when there is no density). On a drillhole layer, interval
  length above each cutoff instead, length-weighted. \cref{ch:gt,sec:gt-holes}.

\item[group] A folder in the layers panel. Its eye hides its children without changing their own
  eyes. A name with slashes files a layer into a group path. \cref{ch:start}.

\item[guard] A routine step, \code{guard: "no-overwrites"}, after which the routine refuses to
  start if a step would replace something that exists. \cref{ch:routines}.

\item[hard / soft boundary] A domain boundary across which grade jumps (hard: estimate each domain
  with its own samples only) or changes gradually (soft: let the estimate use samples from both
  sides near it). Read from a contact plot: a step at the contact with flat sides, or a ramp.
  \cref{ch:contact}.

\item[hole filter] A filter written over collar columns. Its verdict applies to whole holes, across
  every interval table and the traces. \cref{ch:drillholes}.

\item[host] A web page that loads micro in a frame and drives it through \code{window.micro}.
  \cref{app:host}.

\item[inline step] A routine step that carries its tool and settings itself (\code{tool:} with
  \code{with:}) instead of naming a saved recipe. \cref{ch:routines}.

\item[inputs] The questions a recipe or routine asks before it runs, declared under
  \code{inputs:}. Their answers are substituted where \code{\$name} appears. \cref{ch:routines}.

\item[isolate] What a filter does to the view: matching records stay, everything else is neither
  drawn nor pickable. \cref{ch:filter}.

\item[key] The columns that identify a row of a tool's table: series and cutoff for
  grade–tonnage. The key names a failing row and lines rows up between two runs. \cref{ch:checks}.

\item[key join] Bringing columns from one table to another by matching a key column.
  \cref{ch:joins}.

\item[knife] The tool that makes a vertical section along a line you drag in the view. \key{k}.
  \cref{ch:sections}.

\item[lattice] The regular grid a block model's centroids sit on: an origin, a block size and a
  count per axis, and a rotation for a rotated model. \cref{ch:blocks}.

\item[layer] One entry in the layers panel: a dataset as it is shown, with its own eye, color,
  filter and cut mode. \cref{ch:start}.

\item[length-weighted] Of a drillhole statistic: each interval counts in proportion to its length
  (TO − FROM), times its density when one is given, so a 1 m sample weighs a fifth of a 5 m one.
  \cref{sec:gt-holes,ch:resources}.

\item[library] Recipes and routines kept in the browser, outside any project, and runnable without
  one. \cref{ch:library}.

\item[lineage] The derivation tree of a derived layer or export: the operation, its settings, its
  sources and when. \cref{ch:records}.

\item[location] A space of records within an element. A block model has cells; a grid, nodes; a
  mesh, vertices and faces; a drillhole set, collars, the vertices of the hole paths, and intervals. A column
  belongs to one location. \cref{ch:micro}.

\item[machine table] A tool's table as a routine files it: tab-separated, plain column names,
  numbers as they are, an empty cell for a blank. \cref{ch:records}.

\item[materialize] To compute a ƒ column's values once and keep them, as a stored column
  (\ui{Materialize calc column}). \cref{ch:filter}.

\item[materialized column] See \emph{stored column}.

\item[mined volume] The blocks mined in a period, as the reconciliation factors count them: named
  by a period column, inside a solid or above a pit surface (by fraction, nested for pit phases),
  or kept by a filter. \cref{sec:factors-periods}.

\item[never resident] micro's way of reading large files: in passes, a piece at a time, without
  ever holding the whole file decoded in memory. \cref{ch:micro}.

\item[no data] How blanks are drawn: gray, or hidden; never as the lowest value of the scale.
  \cref{ch:color}.

\item[not evaluated] The outcome of a check that read a blank: neither a pass nor a fail. At level
  \code{stop} it stops the run. \cref{ch:checks}.

\item[operation (op)] A recorded, re-runnable derivation: the expression, method or inputs that
  made a column, stored as data in the element manifest. \cref{ch:store}.

\item[pack] A recipe or routine with everything it needs to run: the recipes its steps name and the
  parameter tables it reads. Kept in the library, or exported as one YAML or zip file. Not the same
  thing as a project pack. \cref{ch:library}.

\item[painted column] A category column filled by hand with the brush, or by a classification
  against a solid or a surface. Shown with a ✎. \cref{ch:picking}.

\item[pairing] Matching each sample with the block it falls in, to compare the two values as a
  scatter or a Q-Q plot. \cref{ch:validation}.

\item[parameter table] The table an \code{each} block of a routine fans out over: one run per row.
  Written in the routine, or read from a CSV or a spreadsheet sheet. \cref{ch:routines}.

\item[pick] Clicking a record to see it: it is highlighted, and the record panel shows its full
  row. \cref{ch:picking}.

\item[pinned column] A column held in memory after a scan, within a budget, so that repeated
  filters read no file. \cref{ch:store}.

\item[pre-flight] The checks micro makes on a routine before its first step. If anything is wrong,
  nothing runs. \cref{ch:routines}.

\item[project] A plain folder holding \file{project.json}, the data files by kind, and what micro
  derived from them. \cref{ch:projects}.

\item[project pack] A whole project as one \file{.zip.micro} file: a backup that also opens
  read-only without being unpacked. \cref{ch:store}.

\item[recipe] An analysis setup saved as YAML: the tool and its settings, re-runnable and editable
  by hand. \cref{ch:recipes}.

\item[reconcile] Comparing two models on a common lattice: a map of the difference, and its
  summary. \cref{ch:joins}.

\item[reconciliation factors] The totals of the resource model, grade control and the mill
  compared inside the mined volume, period by period: F1, F2 and F3. The routine tool is
  \code{factors}. \cref{ch:factors}.

\item[record] One row of a location: a block, a point, an interval, a mesh vertex. Its number is
  its position in the data. \cref{ch:micro}.

\item[record panel] The panel that shows a picked record's values, source columns and derived ones
  alike. \cref{ch:picking}.

\item[reinterpret] Opening a layer's source again as another kind of thing: a block model as
  points, a grid as a surface, a table as geometry. \cref{ch:opening}.

\item[relief / flat sheet] The two shapes of a grid shown as a surface: relief takes its height
  from the grid's values; a flat sheet lies at one level and is colored by them. \cref{ch:grids}.

\item[resource table] Tonnage, grade and metal by group and by grade bin or cutoff: the cross-tab
  window, and the table the composer places on a page. Run as the \code{xtab} tool on a drillhole
  layer, it gives length-weighted length, grade and accumulation instead. \cref{ch:resources}.

\item[review layout] The arrangement micro switches to on a narrow screen: one top bar, a bottom bar
  of large targets, and every panel and window as a sheet. \cref{ch:sealed}.

\item[routine] A recipe that runs other steps in order: recipes, inline tools, checks and
  comparisons, with parameter tables. Its whole run can be listed before it starts.
  \cref{ch:routines}.

\item[rule-based styling] Coloring by an ordered list of rules, each an expression with a name and
  a color. A record wears the color of the first rule it matches. \cref{ch:color}.

\item[run folder] The folder a routine run has to itself, \file{runs/<date> <time> <routine>/}:
  its tables, its other outputs and its run record. \cref{ch:records}.

\item[run record] The \file{run.json} a routine run leaves in its run folder: the hash of every
  input and output, and each step's settings, values and checks. Never edited after the run.
  \cref{ch:records}.

\item[run tracker] The window a running routine opens: every expanded step with its live mark, and
  the summary of the run when it ends. Also \emph{the tracker}. \cref{ch:routines}.

\item[scan] One pass over a table that fills each column's count, mean, quantiles and histogram.
  \cref{ch:columns}.

\item[scene] A saved working state: the camera and section, plus each layer's visibility, color
  and filter. \cref{ch:sections}.

\item[scope] Which records an analysis reads: all, the filter's matches, or the 3D selection.
  \cref{ch:analysis}.

\item[Sealed] micro's security posture: the page cannot reach the network, your data stays on your
  device, and it contains no WebAssembly, only JavaScript that can be read. \cref{ch:sealed}.

\item[section] The live cut through the scene: a plane, with a thickness. Blocks and intervals are
  cut exactly at the plane. \cref{ch:sections}.

\item[select through] The selection mode that takes everything inside the shape you drew, hidden or
  not, instead of only what you can see. \cref{ch:picking}.

\item[selection] The set of records chosen with the rectangle or the lasso (or by a tool), per
  layer. It can be painted, exported, used as an analysis scope, or turned into a column.
  \cref{ch:picking}.

\item[series] One line of an analysis: a layer, a column, and optionally a weight and a filter of
  its own. A window compares several. \cref{ch:analysis}.

\item[set manifest] The one element manifest of a drillhole set, \file{drillholes/<set>.element.json}:
  its collar, survey and interval tables, how they relate by hole id, and the desurvey settings.
  Older projects kept the same in a \file{<set>.holes.json}, which still opens.
  \cref{ch:projects,ch:store}.

\item[sheet] In the review layout, a panel or window that rises from the bottom of the screen. One
  is open at a time. \cref{ch:sealed}.

\item[sidecar] A small file micro keeps beside one of yours in a project: a cache
  (\file{.meta.json}), stored columns (\file{.cols/}), a grid's display cache. \cref{ch:store}.

\item[signed distance] A record's distance to the nearest point of a mesh, negative on one side:
  inside a closed solid (decided by winding number), or below an open surface (decided by its
  normals). What \ui{Distance to a surface…} bins, and what its \ui{⇥ column} keeps as
  \col{DIST\_<SURFACE>}. \cref{sec:contact-surface}.

\item[slab] The band of the section that is kept: the plane, plus and minus a thickness.
  \cref{ch:sections}.

\item[solid] A closed mesh, with an inside. micro classifies records against it by winding number.
  An open mesh is a surface, with an above and a below. \cref{ch:geometry}.

\item[spatial join] Resampling one model's columns onto another model's lattice, weighted by
  overlapping volume. \cref{ch:joins}.

\item[statistics scope] What a composer caption's statistics count: the records in the view's
  section, or the whole model. See \emph{captions: in the section}. \cref{sec:compo-statsscope}.

\item[status line] The line along the bottom of the window where micro says, in a sentence, what
  it last did or found, and what could not be done. \cref{ch:start}.

\item[step] One item of a routine's list: a recipe, an inline tool, an \code{each} block, a
  \code{params} sheet, a check, a compare or a guard. \cref{ch:routines}.

\item[stored column] A derived column whose values are kept, one Parquet file per column in the
  data file's \file{.cols/} folder: a ƒ column that was materialized, a join, a broadcast. The
  record panel and the export dialog mark it ✎, numbers and text alike, as they mark painted
  columns; only a fraction inside a solid is marked \% instead. Also called a materialized column. \cref{ch:filter,ch:store}.

\item[sub-blocked model] A block model whose blocks are split into smaller boxes near boundaries.
  Each block keeps its own size and is drawn at it. \cref{ch:blocks}.

\item[swath] Mean grade per band along a direction: the drift plot that compares a model with
  another model or with its samples. \cref{ch:swath}.

\item[table layer] A layer with rows and columns and no geometry: a lookup, a parameter table, a
  kept result. \cref{ch:columns}.

\item[tool] One analysis as a routine can run it: a name, a set of settings, a table and values
  out. Each analysis window is a tool with a window on it. \cref{ch:analysis,app:tools}.

\item[trace] Three things, by context. A sectioned mesh's outline on the cut face
  (\cref{ch:sections}); a hole's path from collar to end (\cref{ch:drillholes}); a fracture trace
  drawn on an outcrop mesh (\cref{ch:structure}).

\item[values] The named numbers a tool returns beside its table, such as a validation's
  \code{bias\_pct}. A check reads them as \code{step.name}. \cref{ch:checks}.

\item[vertical exaggeration] A stretch of the z axis applied to the display only. Coordinates,
  measurements and queries stay real. \cref{ch:sections}.

\item[view (of a model)] The same records shown again as another layer, with its own color,
  filter, visibility and position on the screen. \cref{ch:sections}.

\item[view slot] One of nine numbered bookmarks: \key{g} then a digit goes to it, \key{g} then
  \key{⇧}+digit saves into it. \cref{ch:sections,app:keys}.

\item[watertight] Of a mesh: every edge shared by exactly two triangles. A closed solid, with a
  volume. \cref{ch:qa}.

\item[whole holes] A cut mode for drillholes: a hole whose path crosses the slab is drawn entire,
  and the rest are hidden. \cref{ch:sections}.

\item[zebra] The swath's bands striped onto the 3D view, so the plot and the model can be read
  together. \cref{ch:swath}.

\end{description}
